%!TEX root = ./../../thesis-trier.tex

\paragraph{Sparse Languages.} 
Hoffmann~\cite{DBLP:conf/ictcs/Hoffman20} (inspired 
by a result from~\cite[Thm.~19]{FernauGHHVW19}) 
called a PDFA $\automat{B}{P}{\Sigma}{\mu}{p_0}{F}$  \emph{polycyclic} \IN{polycyclic PDFA}
if, for each $p \in P$, we have
$L(\mathcal B_{p,\{p\}}) \subseteq u_p^*$ for some $u_p \in \Sigma^*$,
and proved the following:
\begin{enumerate}
\item constraint languages accepted by polycyclic automata give
 constrained problems in \NP\ (generalizing a statement from~\cite[Thm. 19]{FernauGHHVW19}),
  		
\item polycyclic automata are equivalently characterized by the property
 that every strongly connected components forms a single cycle,
\item languages accepted by polycyclic PDFAs are closed for union, concatenation
 and quotients,
 
 \IN{nondeterministic automata}
\item nondeterministic automata\footnote{See~\cite{HopUll79} for a definition of nondeterministic automata.} 
 with the same definitional condition accept the same class of languages.
\end{enumerate}
Gawrychowski et al.~\cite[Lem.~2]{DBLP:journals/ijfcs/GawrychowskiKRS10} 
proved that nondeterministic automata fulfilling
the definition condition for polycyclic automata
accept precisely the sparse regular languages.\footnote{The paper by Gawrychowski et al.~\cite{DBLP:journals/ijfcs/GawrychowskiKRS10} 
appeared earlier than~\cite{DBLP:conf/ictcs/Hoffman20},
but the author of~\cite{DBLP:conf/ictcs/Hoffman20} was unaware of the connection to sparse languages
when writing his paper.}  
So, polycyclic automata characterize sparse regular languages (this also follows
from our method of proof used for \cref{constrsync:thm:reg-sparse-or-exp}). 
The characterization of sparse languages
by PDFAs whose strongly connected components form single cycles
was stated by Pin in~\cite[Sec.~XII.4.2]{Pin2022} in terms of forbidden patterns. It also appears
in~\cite[Lem.~5.4]{DBLP:conf/stacs/GanardiHKLM18}.

Also an \NP-hardness condition was given in~\cite[Prop.~8]{DBLP:conf/ictcs/Hoffman20}
that was generalized in \cref{constrsync:thm:NP-hardness}. With this condition it is easy
to find languages accepted by polycyclic constraint PDFAs that give \NP-hard
constrained synchronization problems.

Hoffmann~\cite{DBLP:conf/fct/Hoffmann21} gave a classification for letter-bounded constraint languages
into \NP-complete and polynomial-time solvable constrained problems.
The classification of the present work is a refinement.
Also, the method of proof is different from~\cite{DBLP:conf/fct/Hoffmann21} (for the interested reader 
the original proofs are contained in the arXiv version to which a link is provided in the bibliographic entry~\cite{DBLP:conf/fct/Hoffmann21}, but actually the present
proofs are much more concise). 
A classification for all sparse languages is an \emph{open problem}.
\begin{center}
	\fbox{ \begin{minipage}{30em}
			Open Problem: Classify the computational complexity of \constrsync{L}
			for sparse regular languages $L \subseteq \Sigma^*$.
		\end{minipage}
	}
\end{center}

\paragraph{Monoid Languages and Strongly Connected Constraint Automata.} 
The following results carry over from the unconstrained setting to monoid constraint languages: the pair condition (\cref{constr_sync:thm:pair_condition}); a cubic bound on the length of the shortest synchronizing words in the constraint language (\cref{constr_sync:prop:pair_bound}); a polynomial-time algorithm for computing a synchronizing word of at most cubic length (\cref{constr_sync:thm:monoid_compute_sync_word}); and membership in \NL\ (\cref{thm:reg_monoid}).
No condition is known to identify \NL-hardness for either monoid languages or strongly connected constraint PDFAs; we therefore pose this as an \emph{open problem}.

\begin{center}
	\fbox{ \begin{minipage}{30em}
			Open Problem: Identify the \NL-hard constrained synchronization problem instances 
			for constraint languages that are monoids
			or are given by a strongly connected constraint automaton.
			Furthermore, does a monoid contraint language of exponential growth
			gives an \NL-hard constrained problems?
		\end{minipage}
	}
\end{center}

\IN{returning automaton}
An automaton $\automat{B}{P}{\Sigma}{\mu}{p_0}{F}$ is \emph{returning} if for every state $p \in P$ 
there exists $u \in \Sigma^*$ such that $\applytrans{\mu}{p}{u} = p_0$.
For returning automata $\mathcal B$, in~\cite[Lem. 16 \& Thm.~17]{FernauGHHVW19} the pair condition (\cref{constr_sync:thm:pair_condition})
was established and polynomial-time solvability of $\constrsync{L(\mathcal B)}$ derived.
An accessible automaton
is returning if and only if it is strongly connected. Hence this result is
covered by \cref{constrsync:cor:SCC} of the present work.
In fact, it shows that this is a consequence of the fact that $\constrsync{L}$
is in \NL\ for regular languages $L$ that are monoids.

The languages $(a+b)^*$ and $c(a+b)^*$ from Theorem~\ref{constr_sync:thm:introductory_examples} 
show that concatenating constraint languages might 
increase the computational complexity of the constrained synchronization problem 
w.r.t.~the resulting language
considerably. 
\Cref{constr_sync:subsec:concat} gave conditions for which the complexity never drops when concatenating languages,
which we used to transfer hardness results from known languages. 
On the other hand, the UVW-Condition (Corollary~\ref{constr_sync:cor:UVW}) 
involved the concatenation of three languages 
and yielded constrained problems in \NL.
For example, it implies that $\constrsync{\Sigma^*L}$ 
is in \NL\ for any regular language $L \subseteq \Sigma^*$,
or more generally $\constrsync{IL}$ for any regular ideal $I \subseteq \Sigma^*$ and regular $L \subseteq \Sigma^*$.
A weaker form was stated in~\cite[Thm.~4.1]{DBLP:conf/cocoon/Hoffmann21}, only giving
containment in \PTIME and stating the result for 
PDFAs whose initial state equals their unique final state, a subclass of regular languages that are monoids (see \cref{constr_sync:thm:regular_prefix_codes_and_automata}).


\paragraph{Small Constraint Automata.} Fernau et al.~\cite{FernauGHHVW19}
gave the following classification for small constraint automata with up 
to two states over an at most ternary alphabet.

\begin{theorem}[Fernau et al.~\cite{FernauGHHVW19}]
	\label{thm:classification_MFCS_paper}
	Let $\mathcal B = (\Sigma, P, \mu, p_0, F)$
	be a PDFA.
	If $|P|\le 1$ or $|P| = 2$ and $|\Sigma|\le 2$, then $\constrsync{L(\mathcal B)} \in \PTIME$.
	For $|P| = 2$ with $|\Sigma| = 3$, up to symmetry by renaming of the letters,
	$\constrsync{L(\mathcal B)}$ % $L\mbox{\normalfont \textsc{-Constr-Sync}}$, $L(\mathcal B)\textsc{-Constr-Sync}$
	is $\PSPACE$-complete precisely in the following cases for $L(\mathcal B)$:
	\[ 
	\begin{array}{llll}
		a(b+c)^*        & (a+b+c)(a+b)^*  & (a+b)(a+c)^* & (a+b)^*c \\
		(a+b)^*ca^*     & (a+b)^*c(a+b)^* & (a+b)^*cc^*  & a^*b(a+c)^* \\
		a^*(b+c)(a+b)^* & a^*b(b+c)^*     & (a+b)^*c(b+c)^* & a^*(b+c)(b+c)^*
	\end{array}
	\]
	and polynomial time solvable in all other cases.
\end{theorem}

Theorem~\ref{thm:two-state-classification} extended this result in the $2$-state case 
to arbitrary alphabets and refined the complexity of the polynomial time cases into the three classes~$\AC^0$,
$\LOGSPACE$-complete, and \NL-complete for non-strongly connected constraint PDFAs. 
The extension to arbitrary alphabets for $2$-state constraint PDFAs was
first given in~\cite{DBLP:conf/cocoon/Hoffmann21} (by only distinguishing between \PSPACE-completeness
and~\PTIME). The classification presented here was given in~\cite{HoffmannCocoonExtended},
the extended journal version of~\cite{DBLP:conf/cocoon/Hoffmann21}. 
In~\cite{FernauGHHVW19} the following two 
constrained problems and their complexities given by $3$-state binary PDFAs were presented: \constrsync{ab^*a} (\NP-complete)
and \constrsync{b(aa+ba)^*} (\PSPACE-complete). 
So, the smallest automata giving \NP-complete problems have three states and are over a binary alphabet. 
We completed the picture by giving
computational complexities for all $3$-state binary PDFAs (up to renaming of symbols and states).
We gave a full classification for non-strongly-connected PDFAs
into problems that are in $\AC^0$, \LOGSPACE-, \NL-, \NP-, or \PSPACE-complete.
This classification was also given in~\cite{HoffmannCocoonExtended}.
As a special case of the open problem posed for strongly connected automata
before, we also do not know when the problem
for $3$-state strongly connected binary automata %or strongly connected $2$-state automata 
is \NL-hard.

\begin{center}
	\fbox{ \begin{minipage}{30em}
			Open Problem (special case of the open problem for general strongly connected automata): 
			Identify the \NL-hard constrained synchronization problem instances 
			for $3$-state strongly connected binary constraint PDFAs.
	%		and $2$-state strongly connected constraint PDFAs.
		\end{minipage}
	}
\end{center}

\paragraph{Commutative Languages.} A coarser classification
for commutative regular constraint languages
distinguishing between \PSPACE-complete, \NP-complete
or polynomial-time solvable constrained problems
was given in~~\cite{DBLP:journals/tcs/Hoffmann21}.
Here, the classification was refined
by  subdividing the case
when the constrained problem is in \PTIME into the complexities \NL-complete, \LOGSPACE-complete
and containment in~$\AC^0$.
Furthermore, using the 
general results from \cref{constr_sync:chap:gen_results}
and the cropped Parikh map, the proof was much simplified
and shortened. This simplified proof was also presented in~\cite{HoffmannCocoonExtended}.

\paragraph{Deciding the Complexity.} It is readily seen that 
the classification results for letter-bounded and commutative constraint
languages can be used to decide in polynomial time the computational complexity
of $\constrsync{L}$ when $L \subseteq \Sigma^*$ is given as input (for example, via accepting PDFAs). 

\begin{theorem}
	Let $\mathcal B$ be a PDFA.
	\begin{enumerate}
		\item If $L(\mathcal B)$ is letter-bounded, then we can decide in polynomial time the complexity of \constrsync{L(\mathcal B)},
		i.e., whether it is in $\AC^0$, \LOGSPACE-complete or \NP-complete.
		
		\item If $L(\mathcal B)$ is commutative, then we can decide  in polynomial time the complexity of \constrsync{L(\mathcal B)}, i.e.,
		whether it is in $\AC^0$, \LOGSPACE-complete, \NL-complete, \NP-complete or \PSPACE-complete.
	\end{enumerate}	
\end{theorem} 

Also, for small constraint automata, the complexity up to the fidelity as stated
in the present work can be decided in polynomial time by trimming and comparing
an input PDFA with finitely many alternatives.

\paragraph{A Conjecture for the General Case.} In~\cite[Conclusions and Prospects]{FernauGHHVW19} 
it was asked whether there exist constrained problems complete
for levels in the polynomial time hierarchy other than \NP\ or \PSPACE.
However, we do not believe so and in light of the present results
we formulate the following conjecture.

\begin{center}
	\fbox{ \begin{minipage}{30em}
			Conjecture: For every $L \subseteq \Sigma^*$ the problem \constrsync{L} 
			is either in~$\AC^0$ or complete for one of the following classes: \LOGSPACE, \NL, \NP, \PSPACE. 
		\end{minipage}
	}
\end{center}

In~\cite[Sec.~6]{DBLP:conf/cocoon/Hoffmann21} it was also suspected that only
constrained problems in \PTIME, or complete for \PSPACE\  or \NP\  are possible.

Since every alphabet can be encoded by a solid code over a binary alphabet (as, for example, those given
in \cref{constrsync:ex:binarysolidcode}) which preserves the complexity
by \cref{thm:constr_sync_hom_strongly_self_sync},
confirmation of the conjecture for binary automata implies it for all automata.

In~\cite[Rem.~30]{FernauGHHVW19} it was remarked that the following construction
from \cite[pp. 220-221]{DBLP:journals/mst/Berlinkov14}
can be used to lift hardness results: 
for $\semiautomat{A}{Q}{\{a_1, \ldots, a_k\}}{\delta}$
let $\semiautomat{A}{Q \times \{1,\ldots,k\}}{\{a,b\}}{\delta'}$
with $\delta'((q, i), a) = (q, \min\{i+1,k\})$
and $\delta'((q, i), b) = (\delta(q, a_i), 1)$.
Set $\varphi(a_i) = a^{i - 1}b$ for $i \in \{1,\ldots,k\}$.
Then $u_1 \cdots u_m$ with $u_1, \ldots, u_n \in \{a_1, \ldots, a_k\}$
synchronizes $\mathcal A$ if and only if $b\varphi(u_1) \cdots \varphi(u_m)$
synchronizes $\mathcal A'$. However, since the letter~$b$ has to be read
first this must be allowed by the binary constrained language
restricting the constrained synchronization problem for~$\mathcal A'$ to obtain a reduction.

\paragraph{Road-Coloring and Generalizations.} 
Checking, for a given automaton and word, whether the word synchronizes the automaton
is \LOGSPACE-complete by \cref{prelim:prop:finite_list_sync_word-hardness}.
In~\cite{DBLP:journals/jcss/VorelR19} the authors classified the computational complexity of the so-called \emph{road-coloring variant of synchronization}: 
given a digraph of constant outdegree and a word $w$, whether the digraph
admits a labeling of its edges with letters such that a synchronizing automaton with synchronizing
word~$w$ results. Depending on the class of admissible input words, the problem
is either polynomial-time solvable or \NP-complete.
Now this suggests the following generalization of $\constrsync{L}$: for a fixed constraint language,
given a digraph of constant outdegree, does it admit a labeling of its edges such
that a synchronizing automaton results having a synchronizing word in the constraint language.
%Note that here, as in $\constrsync{L}$, we have fixed the constraint language.
In~\cite{DBLP:journals/jcss/VorelR19} it has been proved that this problem is \NP-complete
even for the constraint language $\{abb\}$.
When we admit the constraint language to be part of the input,
for $L = \Sigma^n$, where~$n$ is an input parameter, this problem has been investigated
in~\cite{ROMAN2015383,DBLP:journals/ipl/Roman11}. The generalization
for arbitrary languages as input has been mentioned 
in~\cite{DBLP:journals/ipl/Roman11,DBLP:journals/jcss/VorelR19}.
Other generalizations where the constraint language is part of the input and belongs
to some class of admissible constraint languages for our constrained problem or the road-coloring
variant are conceivable, but we have not pursued to investigate these variants here.